\Askhsh{36101}{2}
\begin{ylh}{1}
    \item 4.4. Γωνίες με πλευρές παράλληλες
    \item 4.6. Άθροισμα γωνιών τριγώνου.
\end{ylh}
% ===================================================================
%  Αρχείο: 36101.tex (Fragment)
%  Αυτόματη μετατροπή μέσω doc2latex.py
% ===================================================================

ΘΕΜΑ 2

Δίνεται τρίγωνο ΑΒΓ με $\widehat{A} = 80{^\circ}$ και $\widehat{B} = 20{^\circ} + \widehat{\Gamma}$, και ΑΔ η διχοτόμος της γωνίας $\widehat{A}$.

\begin{alist}
\item Να υπολογίσετε τις γωνίες $\widehat{B}$ και $\widehat{\Gamma}$. \monades{12}

\item Φέρουμε από το $\Delta$ ευθεία παράλληλη στην $AB$, που τέμνει την $A\Gamma$ στο $E$. Να υπολογίσετε τις γωνίες $A\widehat{\Delta}E$ και $E\widehat{\Delta}\Gamma$. \monades{13}

\begin{center}
\begin{tikzpicture}[scale=1]
\tkzDefPoint(0,0){B}
\tkzDefPoint(5,0){C}
\tkzDefTriangle[two angles = 60 and 40](B,C) \tkzGetPoint{A}

\tkzDefLine[bisector](B,A,C) \tkzGetPoint{AD_dir}
\tkzInterLL(A,AD_dir)(B,C) \tkzGetPoint{D}

\tkzDefLine[parallel=through D](A,B) \tkzGetPoint{DE_dir}
\tkzInterLL(D,DE_dir)(A,C) \tkzGetPoint{E}

\draw[l3] (A)--(B)--(C)--cycle;
\draw[l3] (A)--(D);
\draw[l3] (D)--(E);

\tkzDrawPoints(A,B,C,D,E)
\tkzLabelPoint[above](A){$A$}
\tkzLabelPoint[below](B){$B$}
\tkzLabelPoint[below](C){$\Gamma$}
\tkzLabelPoint[below](D){$\Delta$}
\tkzLabelPoint[right](E){$E$}

\tkzMarkAngle[size=0.6,mark=none](B,A,D)
\tkzMarkAngle[size=0.6,mark=none](D,A,C)
\tkzMarkAngle[size=0.5,mark=none](C,D,E)
\end{tikzpicture}
\end{center}
\end{alist}
